\section{Conclusion}
\label{Conclusion}
%It outlines how the publication has answered the research questions or the research hypotheses. For this, research questions or the research hypotheses must explicitly (by using their names) be mentioned. Here, it is crucial to outline very explicitly what is new.


A comprehensive approach that integrates ML, CTs and experiments is presented. This solves the lack of transferable and interpretable materials laws for complex composite materials and assemblies. The surrogate design model provides an efficient way to overcome the traditional bottleneck in the design-test cycle. 

The algorithmic representation ensures the interpretability of the neural network-based models. The extraction of explicit, formula-based models enhances understanding and facilitates meaningful contributions to materials design. 

In addition, the morphological representation provides a visual understanding of the materials system through synthesis of 3D tomography data. This enables the analysis of manufacturing influences on morphological structures and provides valuable insights into the material design process.

The proposed framework not only streamlines the materials design process but also empowers researchers with a powerful tool to explore and predict the behavior of complex material composites. By combining CTs synthesis and explainable AI, this modern approach opens opportunities for the rapid and informed design of materials with tailored properties.